\documentclass[12pt,a4paper]{article}
\usepackage{ugmath}
\usepackage{float}
\usepackage{placeins}
\usepackage{array}
\usepackage[labelfont=bf,labelsep=space]{caption}
\newcommand{\paren}[1]{\left( #1 \right)}

\title{Solucion de Problemas\\
Taller de calculo}

\author{Ricardo Leon Martinez}

\begin{document}
    \maketitle

    \section*{Problema 1}

    Sea $f:(0,\infty)\to\R$ como
    \begin{equation*}
        f(x)=2^{-x}+2^{-\frac{1}{x}}.
    \end{equation*}

    \begin{enumerate}
        \item[(a)] Calcula $\lim_{x\to0^{+}}f(x)$ y $\lim_{x\to\infty}f(x)$.
        \item[(b)] Calcula (y demuestra que tu cálculo es correcto) el máximo de $f$ en
        todo su dominio.
        \item[(c)] ¿Qué puedes decir de su mínimo (también global)? 
    \end{enumerate}

    \begin{solucion}
        \subsection*{(a)}
        Calcularemos primero $\lim_{x\to0^{+}}f(x)$. Observemos primero que
        \begin{equation*}
            2^{-x}=e^{-x\ln(2)}.
        \end{equation*}
        Como $x\to0^{+}$, se sigue que
        \begin{equation*}
            -x\ln(2)\to0,
        \end{equation*}
        y por continuidad de la exponencial,
        \begin{equation*}
            2^{-x}\to e^{0}=1.
        \end{equation*}
        Por otro lado,
        \begin{equation*}
            2^{-\frac{1}{x}}=e^{-\frac{\ln(2)}{x}}.
        \end{equation*}
        Dado que $\frac{1}{x}\to+\infty\quad(x\to0^{+})$, tenemos
        \begin{equation*}
            -\frac{\ln(2)}{x}\to\infty,
        \end{equation*}
        de donde $2^{-\frac{1}{x}}\to0$. Por consiguiente
        \begin{equation*}
            \lim_{x\to0^{+}}f(x)=1.
        \end{equation*}
        Ahora calculamos
        \begin{equation*}
            \lim_{x\to\infty}f(x).
        \end{equation*}
        Cuando $x\to\infty$, $-x\ln(2)\to-\infty$, por lo que $2^{-x}\to0$. Además,
        \begin{equation*}
            \frac{1}{x}\to0,
        \end{equation*}
        de modo que
        \begin{equation*}
            -\frac{\ln(2)}{x}\to0.
        \end{equation*}
        y entonces
        \begin{equation*}
            2^{-\frac{1}{x}}=1.
        \end{equation*}
        Por tanto,
        \begin{equation*}
            \lim_{x\to\infty}f(x)=1.
        \end{equation*}

        \subsection*{(b)}

        Derivamos
        \begin{equation*}
            f^{\prime}(x)-(\ln(2))2^{-x}+\frac{\ln(2)}{x^{2}}2^{-\frac{1}{x}}.
        \end{equation*}
        Factorizando $\ln(2)\gt0$,
        \begin{equation*}
            f^{\prime}(x)=\ln(2)\paren{\frac{2^{-1/x}}{x^{2}}-2^{-x}}.
        \end{equation*}
        Por lo tanto,
        \begin{equation*}
            f^{\prime}(x)=0\Longleftrightarrow\frac{2^{-1/x}}{x^{2}}=2^{-x}.
        \end{equation*}
        Equivalentemente,
        \begin{equation*}
            2^{x-1/x}=x^{2}.
        \end{equation*}
        Tomando logaritmo en base 2,
        \begin{equation*}
            x-\frac{1}{x}=2\log_{2}(x).
        \end{equation*}
        Es inmediato ver que $x=1$ es la solución. Veamos ahora que es única.

        Definamos
        \begin{equation*}
            g(x)=x-\frac{1}{x}-2\log_{2}(x).
        \end{equation*}
        Entonces
        \begin{equation*}
            g^{\prime}(x)=1+\frac{1}{x^{2}}-\frac{2}{x\ln(2)}.
        \end{equation*}
        Además,
        \begin{equation*}
            g^{\prime\prime}(x)=-\frac{2}{x^{3}}+\frac{2}{x^{2}\ln(2)}=\frac{2}{x^{3}}\paren{\frac{x}{\ln(2)}-1}.
        \end{equation*}
        En particular, $g^{\prime\prime}$ cambia de signo una sola vez, de modo que $g^{\prime}$
        tiene un único mínimo. Como
        \begin{equation*}
            g^{\prime}(1)=2-\frac{2}{\ln(2)}\lt0,
        \end{equation*}
        y
        \begin{equation*}
            g^{\prime}(x)\to\infty\quad\text{ cuando }\quad x\to0^{+}\text{ o }x\to\infty,
        \end{equation*}
        se concluye que $g^{\prime}$ tiene exactamente dos ceros y, por tanto, $g$ tiene
        un único mínimo global. Ahora,
        \begin{equation*}
            g(1)=0.
        \end{equation*}
        Como el mínimo global de $g$ se alcanza precisamente en $x=1$, se sigue que $g(x)\geq0$
        para todo $x\gt0$, con igualdad únicamente en $x=1$. Por consiguiente,
        \begin{equation*}
            f^{\prime}(x)=
            \begin{cases}
                f(x)\gt0,\quad0\lt x\lt1,\\
                f(x)=0,\quad x=1,\\
                f(x)\lt0,\quad x\gt1.
            \end{cases}
        \end{equation*}
        Así, $f$ es creciente en $(0,1)$ y decreciente $(1,\infty)$. Luego $x=1$ corresponde a un
        máximo global. Finalmente,
        \begin{equation*}
            f(1)=2^{-1}+2^{-1}=1.
        \end{equation*}
        Concluimos que el máximo global de $f$ es 1, y se alcanza únicamente en $x=1$.

        \subsection*{(c)}

        Para todo $x\gt0$,
        \begin{equation*}
            2^{-x}\gt0,\quad 2^{-1/x}\gt0,
        \end{equation*}
        y por tanto $f(x)\gt0$. Además,
        \begin{equation*}
            \lim_{x\to0^{+}}f(x)=1,\qquad\lim_{x\to\infty}f(x)=1.
        \end{equation*}
        Como la función es positiva y alcanza su valor máximo en $x=1$, se tiene
        \begin{equation*}
            0\lt f(x)\leq1.
        \end{equation*}
        Sin embargo, la función nunca toma el valor 0. En consecuencia, $f$ no tiene mínimo global.
    \end{solucion}

    \section*{Problema 2}

    \textit{(a)} Considera la párabola $y=ax^{2}$. Toma sobre su gráfica dos puntos
    cualesquiera $A=(x_{1},y_{1})$ y $B=(x_{2},y_{2})$. Considera la cuarda $\overline{AB}$.
    Prueba que existe un único punto $C$ sobre la gráfica de la parábola, tal que la recta
    tangente a la párabola en $C$ es paralela a la cuerda $\overline{AB}$. Calcula
    las coordenadas de $C$.

    \textit{(b)} Vas a probar el siguiente teorema de Arquímedes: el área de la region
    dentro de la parábola y limitada por la cuerda $\overline{AB}$ es $4/3$ del área del
    triangulo $\Delta ABC$. (no copie lo demas porque se me hizo demasiado).

    \textit{(c)} Explica por qué el teorema es cierto para cualquier parábola y no únicamente
    para $y=ax^{2}$.

    \begin{solucion}
        \subsection*{(a)}
        Sean
        \begin{equation*}
            A=(x_{1},ax_{1}^{2}),\qquad B=(x_{2},ax_{2}^{2}).
        \end{equation*}
        La pendiente de la cuerda $AB$ es
        \begin{equation*}
            m_{AB}=\frac{ax_{2}^{2}-ax_{1}^{2}}{x_{2}-x_{1}}=a(x_{1}+x_{2}).
        \end{equation*}
        Por otra parte, si $f(x)=ax^{2}$, entonces
        \begin{equation*}
            f^{\prime}(x)=2ax.
        \end{equation*}
        Así, la recta tangente en un punto
        \begin{equation*}
            C=(x_{c},ax_{c}^{2})
        \end{equation*}
        tiene pendiente $2ax_{c}$. Queremos que dicha tangente sea paralela a la cuerda $AB$,
        por lo que debe cumplirse
        \begin{equation*}
            2ax_{c}=a(x_{1}+x_{2}).
        \end{equation*}
        Suponiendo $a\neq0$, obtenemos
        \begin{equation*}
            x_{c}=\frac{x_{1}+x_{2}}{2}.
        \end{equation*}
        En consecuencia,
        \begin{equation*}
            C=\paren{\frac{x_{1}+x_{2}}{2},a\paren{\frac{x_{1}+x_{2}}{2}}^{2}}.
        \end{equation*}
        Además, esta elección es única, pues la pendiente de la tangente depende únicamente
        de la coordenada $x$.

        \subsection*{(b)(i)}
        Consideremos ahora las cuerdas $AC$ y $CB$. Sean $D$ y $E$ los puntos de la parábola
        tales que las tangentes en dichos puntos son paralelas, respectivamente, a $AC$ y
        $BC$. Sea ademas $M$ el punto donde la recta vertical que pasa por $C$ intersecta
        a la cuerda $AB$. Observemos que dicha recta bisecta la cuerda $AB$. En efecto,
        en el inciso anterior se obtuvo que
        \begin{equation*}
            x_{C}=\frac{x_{1}+x_{2}}{2},
        \end{equation*}
        de modo que la vertical por $C$ pasa por el punto medio de los extremos de la cuerda.
        El objetivo es demostrar que
        \begin{equation*}
            \ar(\Delta ADC)=\frac{1}{4}\ar(\Delta AMC),
        \end{equation*}
        y
        \begin{equation*}
            \ar(\Delta CEB)=\frac{1}{4}\ar(\Delta MCB).
        \end{equation*}
        ambas demostraciones son análogas, así que probaremos únicamente
        la segunda. Sea $G$ el punto donde la vertical por $E$ intersecta
        a $CB$, y sea $H$ el punto donde dicha vertical intersecta al segmento
        $MB$. Mostraremos primero que la vertical por $E$ bisecta tanto a $CB$ como a $MB$.
        En efecto, si
        \begin{equation*}
            C=(x_{c},ax_{x}^{2}),\qquad B=(x_{2},ax_{2}^{2}),
        \end{equation*}
        entonces, por el inciso (a), el punto $E$ satisface
        \begin{equation*}
            x_{E}=\frac{x_{c}+x_{2}}{2},
        \end{equation*}
        pues la tangente en $E$ es paralela a la cuerda $CB$. Por tanto, la vertical
        por $E$ pasa por el punto medio horizontal de $CB$, y en consecuencia bisecta
        dicho segmento. Análogamente, como $M$ y $B$ tiene coordenadas horizontales
        $x_{c}$ y $x_{2}$, respectivamente, la misma vertical también bisecta el
        segmento $MB$. Ahora veamos que los triangulos $\Delta EGH$ y $\Delta MCB$
        son semejantes. En efecto, ambos son triangulos son rectangulos, comparten la inclinación
        determinada por la cuerda $CB$ y sus lados verticales y horizontales son parelelos
        entre sí. Además, como $G$ es el punto medio de $CB$ y $H$ es el punto medio de $MB$,
        las longitudes horizontales satisfacen
        \begin{equation*}
            GH=\frac{1}{2}MC.
        \end{equation*}
        Por otra parte, usando que $E$ pertenece a la parábola, una cuenta directa muestra que
        \begin{equation*}
            EG=\frac{1}{2}GH
        \end{equation*}
        Por tanto,
        \begin{equation*}
            \frac{EG}{MC}=\frac{1}{4}.
        \end{equation*}
        Ahora bien, los triángulos $\Delta CEB$ y $\Delta MCB$ tiene bases situadas sobre la misma
        recta $CB$, y sus alturas respectivas son $EG$ y $MC$. Luego,
        \begin{equation*}
            \frac{\ar(\Delta CEB)}{\ar(\Delta MCB)}=\frac{EG}{MC}=\frac{1}{4}.
        \end{equation*}
        Así,
        \begin{equation*}
            \ar(\Delta CEB)=\frac{1}{4}\ar(\Delta MCB).
        \end{equation*}
        De manera completamente análoga se obtiene
        \begin{equation*}
            \ar(\Delta ADC)==\frac{1}{4}\ar(\Delta AMC).
        \end{equation*}
        Finalmente,
        \begin{equation*}
            \ar(\Delta ADC)+\ar(\Delta CEB)=\frac{1}{4}(\ar(\Delta AMC)+\ar(\Delta MCB)).
        \end{equation*}
        Como
        \begin{equation*}
            \ar(\Delta AMC)+\ar(\Delta MCB)=\ar(\Delta ACB),
        \end{equation*}
        concluimos que
        \begin{equation*}
            \ar(\Delta ADC)+\ar(\Delta CEB)=\frac{1}{4}\ar(\Delta ACB).
        \end{equation*}

        \subsection*{(b)(ii)}

        El paso anterior puede repetirse indefinidamente. Cada uno de los triángulos obtenidos
        produce nuevas cuerdas y, por tanto, nuevos triángulos inscritos. Ademas, la suma de las
        áreas de los nuevos triangulos es igual a una cuarta parte de la suma de las áreas de
        la generación anterior. Si denotamos por
        \begin{equation*}
            T=\ar(\Delta ABC),
        \end{equation*}
        la suma total de areas obtenidas es
        \begin{equation*}
            T+\frac{T}{4}+\frac{T}{16}+\cdots
        \end{equation*}
        Es una serie geometrica de razón $1/4$, cuya suma es
        \begin{equation*}
            \frac{T}{1-\frac{1}{4}}=\frac{4}{3}T.
        \end{equation*}
        Por consiguiente, el área del segmento parabólico es
        \begin{equation*}
            \frac{4}{3}\ar(\Delta ABC).
        \end{equation*}

        \subsection*{(C)}

        Toda parabola puede obtenerse a partir de la otra mediante traslaciones, rotaciones
        y cambios de escala. Esta transformaciones preservan las relaciones de paralelismo
        y escalan todas las áreas por el mismo factor. En consecuencia, la proporción
        \begin{equation*}
            \frac{4}{3}
        \end{equation*}
        permanece invariante. Por tanto, el teorema es válido para cualquier parábola y no
        unicamente para aquellas de la forma
        \begin{equation*}
            y=ax^{2}.
        \end{equation*}
    \end{solucion}
    

\end{document}
